% Propietario conservado: /Users/ruben/Documents/New project/output/EXTREMA_MEDIA_RAZON_RECIPROCIDAD_ES_EN_20260918/ARTICULO/ES/source/nuclear/05.tex
\chapter{Reconstruction d’incidence et covariance}
\label{x_reciprocidad:nuclear:5}
\section{L’isométrie complète et son domaine}\label{x_reciprocidad:nuc05:isometria}

Soit \(H=\mathbb R^{12}\), muni du produit euclidien ; \(P_0=I-\mathbf1\mathbf1^*/12\), \(\mu(k)=\mathbf1^*k/12\), et \(\mathcal I:H\to\mathbb R^{132}\) la matrice d’incidence de \cref{x_reciprocidad:exc:entrelazador,x_reciprocidad:exc:isometria}. Nous écrivons \(E_{\rm inc}=\mathcal I(\mathbf1^\perp)\) et

\[
Y=\mathbb R\oplus E_{\rm inc},\qquad
\langle(\mu,y),(\nu,z)\rangle_Y=12\mu\nu+\langle y,z\rangle.
\]

L’isométrie complète construite dans la section~\ref{x_reciprocidad:nuclear:2} est

\[
F:H\longrightarrow Y,\qquad Fk=(\mu(k),\mathcal I P_0k/6),
\]

\[
F^{-1}(\mu,y)=\mu\mathbf1+P_0\mathcal I^*y/6.
\]

La matrice de Gram \(\mathcal I^*\mathcal I=36I+30\mathbf1\mathbf1^*\) prouve les deux compositions inverses et \(F^*F=I_H\), \(FF^*=I_Y\), avec le produit pondéré indiqué. C’est pourquoi \(F\) est orthogonal entre \(H\) et \(Y\). Il peut être complexifié sans changer les preuves.

Le codomaine est \(Y\), de dimension douze, et non tout \(\mathbb R\oplus\mathbb R^{132}\). Étendre un opérateur de \(Y\) à son complément dans cet espace ambiant exige un choix supplémentaire ; cela ne fait pas partie de la dynamique transportée de manière unique. En coordonnées signées normalisées, on utilise \(FJ_H\), avec \(J_H=H_{12}/2\). L’inversion entière reste \(H_{12}/4\) dans son sous-réseau, et non un normalisateur métrique alternatif.

\section{Théorème de relèvement covariant dans l’image}\label{x_reciprocidad:nuc05:levantamiento}

\begin{theorem}
Pour un opérateur borné \(D\) sur \(H\), il existe un unique opérateur sur \(Y\) qui satisfait \(D_YF=FD\). C’est

\[
D_Y=FDF^{-1}.
\]

La correspondance conserve les sommes, les produits, les adjoints, la norme, la positivité et l’identité. En particulier, elle envoie les unitaires sur des unitaires et les observables autoadjoints sur des observables autoadjoints. Pour un unitaire \(C\) et un observable \(A\),

\[
C^*AC=A\quad\Longleftrightarrow\quad C_Y^*A_YC_Y=A_Y.
\]

Si \(R\) est un changement unitaire de carte, \(R_Y=FRF^{-1}\), alors


\[
F(RDR^{-1})F^{-1}=R_YD_YR_Y^{-1}.
\]
\end{theorem}

\begin{proof}
La première formule s’obtient en multipliant l’entrelacement par \(F^{-1}\). L’unicité découle de la surjectivité sur \(Y\). Les identités de produits et d’adjoints résultent de \(F^{-1}=F^*\), et l’égalité des normes de l’isométrie surjective. La dernière égalité est une composition associative de ces applications. Aucune hypothèse de commutation entre \(R\) et \(D\) n’intervient.
\end{proof}

Pour \(g\in\operatorname{Aut}(\mathcal H)\), la proposition~\ref{x_reciprocidad:exc:isometria} fournit

\[
F\rho_{12}(g)=\bigl(1\oplus\rho_{\rm hex}(g)\bigr)F.
\]

Ainsi, le changement de carte possède une action combinatoire concrète sur les hexades. En revanche, pour un unitaire général \(D\), \(D_Y\) est un opérateur de la représentation incidencielle, et non nécessairement une permutation du plan. Préserver le produit scalaire de cette représentation ne transforme pas tous ses unitaires en automorphismes combinatoires ni en éléments du Monstre.

\subsection{Covariance sous un changement de carte non commutant}

Soit \(C\) l’avancée \(Ce_p=e_{p+1}\) cyclique sur les douze positions, c’est-à-dire le sens inverse de la convention \(S\) de \cref{x_reciprocidad:sec:k-registro}. C’est une opération de la carte des fenêtres, non l’holonomie temporelle complète. Soit \(R\) l’étoile déjà construite

\[
R=(1\ 3)(2\ 11)(5\ 7)(8\ 12).
\]

Alors \(RCe_1=e_{11}\), tandis que \(CRe_1=e_4\). Ils ne commutent pas. Cependant, pour \(C'=RCR^{-1}\), \(T(C')R=RT(C)\) et \(\eta(C')R=R^{\oplus9}\eta(C)\) sont exactement satisfaits, et de même dans \(Y\) au moyen de \(F\). La covariance de carte fonctionne effectivement dans un exemple où la commutation en carte fixe échoue.


\section{Récupération explicite à partir du registre d’incidence des compléments}\label{x_reciprocidad:nuc05:recuperacion}

Soient \(C\) unitaire, \(T=(8I+C)/9\), \(\eta=(I-JJ^*)S_CJ:H\to H^9\), et \(P\) la projection sur \(\operatorname{Fix}(C)\), comme dans la section~\ref{x_reciprocidad:nuclear:2}. Nous écrivons \(F_9=F^{\oplus9}\). La version incidencielle du registre fini est

\[
\mathcal W_Nk=\left(FT^Nk,F_9\eta k,F_9\eta Tk,\ldots,F_9\eta T^{N-1}k\right).
\]

\begin{theorem}
Cette application est isométrique et possède l’inverse à gauche explicite

\[
\mathcal W_N^*(a,b_0,\ldots,b_{N-1})
=T^{*N}F^{-1}a+
\sum_{j=0}^{N-1}T^{*j}\eta^*F_9^{-1}b_j.
\]

En particulier, appliquer cette formule à \(\mathcal W_Nk\) récupère exactement \(k\). Le registre infini

\[
\mathcal W_\infty k=\left(FPk,(F_9\eta T^jk)_{j\ge0}\right)
\]

est également isométrique, et son inverse à gauche est

\[
k=PF^{-1}a+\sum_{j\ge0}T^{*j}\eta^*F_9^{-1}b_j
\quad\text{pour }(a,(b_j))=\mathcal W_\infty k.
\]

La série converge en norme. Pour des données arbitraires de carré sommable, la même formule définit l’adjoint borné, bien que toute donnée du codomaine n’appartienne pas à l’image de \(\mathcal W_\infty\).
\end{theorem}

\begin{proof}
Appliquer \(F\) et \(F_9\) à chaque composante de \(V_N\) ou \(V_\infty\) conserve le produit scalaire. Prendre l’adjoint de cette composition donne les formules. L’identité télescopique de la section~\ref{x_reciprocidad:nuclear:2} prouve \(\mathcal W_N^*\mathcal W_N=I\), et sa limite forte prouve \(\mathcal W_\infty^*\mathcal W_\infty=I\). La convergence de la série provient de l’application d’un opérateur borné aux troncatures d’une suite dans une somme directe hilbertienne.
\end{proof}

Cette reconstruction ne conserve pas \(k\) comme une première coordonnée intacte : la première composante est le terminal comprimé ou son secteur fixe ; le reste provient des compléments successifs. Lorsque \(P=0\), les compléments incidenciels retrouvent à eux seuls \(k\). La contribution retrouvée à partir des \(N\) premiers blocs est

\[
k_N=\sum_{j<N}T^{*j}\eta^*F_9^{-1}b_j
=\bigl(I-T^{*N}T^N\bigr)k.
\]

Le résidu est \(T^{*N}T^Nk\), qui tend vers zéro si \(P=0\). Avec un secteur fixe non nul, il faut conserver \(FPk\) ; la mémoire des compléments seule ne récupère pas ce secteur. Les taux de reconstruction et la dynamique sur l’indice des blocs sont développés dans la section~\ref{x_reciprocidad:nuclear:3}.

Si \(A\) est auto-adjoint et \([A,C]=0\), la forme de l’observable est également conservée :

\[
\langle k,Ak\rangle
=\langle FPk,A_YFPk\rangle_Y
+\sum_{j\ge0}\langle F_9\eta T^jk,(I_9\otimes A_Y)F_9\eta T^jk\rangle.
\]

La série scalaire converge absolument parce que \(A\) est borné et que la somme des carrés des normes des composantes est finie. L’hypothèse exprime la conservation de \(A\) par \(C\) ; elle n’exige pas la commutation de \(C\) avec les changements de carte exceptionnels.


\section{Relèvement naturel du raffinement avec des repères locaux}\label{x_reciprocidad:nuc05:refinamiento}

La construction de \cref{x_reciprocidad:iv:cont:historias,x_reciprocidad:iv:cont:medida} fournit des ensembles d’états cylindriques enrichis \(Z_n\), des troncatures et des mesures projectives à support total. Pour un descendant \(z'\) de \(z\),

\[
w_n(z'|z)=\sqrt{\mu_{n+1}(z')/\mu_n(z)},\qquad
\sum_{\tau z'=z}w_n(z'|z)^2=1.
\]

L’étiquette \(z'\) conserve le registre ordonné, l’orientation, la retenue, la frontière, la mémoire, la route et le terminal. Le raffinement d’un état n’identifie pas des descendants distincts.

Soient \(q(z)\) le registre entier de mémoire, \(C\) un unitaire sur \(H\), et \(R_z\) des changements orthogonaux de repère de la fibre. Pour des repères d’incidence concrets, on peut utiliser \(R_z=\rho_{12}(g_z)\), avec \(g_z\in\operatorname{Aut}(\mathcal H)\), en transportant leurs marques. Nous définissons

\[
G_z=R_zC^{q(z)},\qquad
U_{z',z}=G_{z'}G_z^*
=R_{z'}C^{q(z')-q(z)}R_z^*.
\]

Cette construction prolonge l’équivalence de jauge explicite de \cref{x_reciprocidad:exc:doble-lectura}. Elle n’exige pas \([R_z,C]=0\) et laisse visible l’ordre des opérations.

Sur \(\mathscr H_n=\ell^2(Z_n)\otimes H\),

\[
\mathscr U_n(e_z\otimes v)
=\sum_{\tau z'=z}w_n(z'|z)e_{z'}\otimes U_{z',z}v.
\]

Nous définissons \(\mathscr F_n=I_{\ell^2(Z_n)}\otimes F\), de codomaine \(\mathscr Y_n=\ell^2(Z_n)\otimes Y\).

\begin{theorem}
Le raffinement est isométrique et satisfait

\[
\mathscr U^Y_n\mathscr F_n=\mathscr F_{n+1}\mathscr U_n,
\quad
\mathscr U^Y_n=\mathscr F_{n+1}\mathscr U_n\mathscr F_n^{-1}.
\]

La formule de chaque terme est la précédente avec \(U^Y_{z',z}=FU_{z',z}F^{-1}\). Les systèmes de plusieurs profondeurs conservent la composition et admettent une isométrie surjective entre leurs limites inductives. Tous les coefficients de fibre sont préservés exactement avec leurs étiquettes \(z\).
\end{theorem}

\begin{proof}
Les descendants de parents distincts sont orthogonaux ; chaque \(U_{z',z}\) est unitaire ; et les poids ont une somme des carrés égale à un. Cela prouve l’isométrie. Dans une trajectoire, les facteurs intermédiaires \(G_z^*G_z\) s’annulent et les quotients de mesures se télescopent. La composition ne dépend pas de la partition de la trajectoire. Le carré avec \(F\) se vérifie sur \(e_z\otimes v\), terme à terme. Les applications et leurs inverses sont des isométries compatibles, de sorte qu’elles s’étendent réciproquement aux complétions inductives.
\end{proof}

\subsection{Observables et dynamique compatibles}

Si \(A_0\) est un observable borné de la fibre, soit

\[
(\mathscr A_n\psi)_z=G_zA_0G_z^*\psi_z.
\]

Alors \(\mathscr A_{n+1}\mathscr U_n=\mathscr U_n\mathscr A_n\), d’où \(\mathscr U_n^*\mathscr A_{n+1}\mathscr U_n=\mathscr A_n\). Cela s’obtient en substituant \(U_{z',z}=G_{z'}G_z^*\). C’est la conservation d’un champ d’observables transporté par raffinement ; cela n’affirme pas que \(A_0\) commute avec \(C\) dans une carte fixe.

La dynamique unitaire par niveau

\[
(\mathscr C_n\psi)_z=G_zCG_z^*\psi_z
\]

satisfait \(\mathscr C_{n+1}\mathscr U_n=\mathscr U_n\mathscr C_n\). Ils entrelacent également leurs adjoints : puisque les deux \(\mathscr C\) sont unitaires, l’identité est multipliée par leurs inverses. Ils entrelacent donc \(T\), leurs compléments nonadiques et leurs projections fixes. Le registre fini et le registre infini de la section~\ref{x_reciprocidad:nuc05:recuperacion} sont naturels par rapport à ce raffinement. En appliquant \(\mathscr F_n\), ces mêmes carrés sont valables dans les fibres incidentielles. Pour que \(\mathscr A_n\) soit en outre conservé par cette dynamique à un niveau fixe, la condition précise est \([A_0,C]=0\) ; la covariance entre repères ne requiert toujours pas \([R_z,C]=0\).

\subsection{Spécialisation de mémoire bilatérale avec repères variables}

Dans l’orbite réversible de \cref{x_reciprocidad:nuclear:3}, une publication de coordonnées peut être représentée par \(\psi=(\psi_m)_{m\in\mathbb Z}\in\ell^2(\mathbb Z;H)\). Son avancée est \((C_0\psi)_m=\psi_{m-1}\). \(F\) agit point par point. Pour des repères \(R_m=\rho_{12}(g_m)\), le changement unitaire \((\mathscr R\psi)_m=R_m\psi_m\) produit

\[
(C'\psi)_m=R_mR_{m-1}^*\psi_{m-1}.
\]

L’opérateur incidenciel a la même formule avec \(1\oplus\rho_{\rm hex}(g_mg_{m-1}^{-1})\). On n’impose ni des repères constants ni la commutation avec le décalage. Comme \(C'\) est conjugué à \(C_0\), il n’a pas de vecteurs fixes non nuls dans cet espace de suites de carré sommable. La section~\ref{x_reciprocidad:nuc05:recuperacion} retrouve la publication linéaire complète à partir de ses compléments incidenciels. Il s’agit de l’orbite et des coordonnées représentées, et non d’une identification de toutes les histoires TPK à \(\mathbb Z\).

\section{Extension directe aux douze fibres \texorpdfstring{\(A_2\)}{A2}}\label{x_reciprocidad:nuc05:a2}

La construction exceptionnelle conserve douze fibres \(A_2\) et l’opérateur \(g_c\) d’ordre trois développé dans la section~\ref{x_reciprocidad:nuc04:c3}. Soit \(V_{A_2}=A_2\otimes_{\mathbb Z}\mathbb R\) sa réalisation euclidienne bidimensionnelle. La même isométrie positionnelle peut être tensorisée sur \(\mathbb R\) avec cet espace :

\[
F_{A_2}=F\otimes I_{V_{A_2}}:
\mathbb R^{12}\otimes_{\mathbb R}V_{A_2}
\longrightarrow Y\otimes_{\mathbb R}V_{A_2}.
\]

Sa première composante est la moyenne vectorielle et sa seconde est constituée des sommes incidentielles centrées, maintenant à valeurs dans \(V_{A_2}\). Le gramien d’incidence, tensorisé avec l’identité, prouve de nouveau l’isométrie complète et l’inverse. C’est pourquoi \(g_c^Y=F_{A_2}g_cF_{A_2}^{-1}\) est orthogonal, d’ordre trois et sans point fixe. La composition nonadique de la section~\ref{x_reciprocidad:nuclear:2} peut se publier et se récupérer sur cet écran à valeurs dans \(V_{A_2}\), au moyen des sections~\ref{x_reciprocidad:nuc05:levantamiento} et~\ref{x_reciprocidad:nuc05:recuperacion}. On ne confond pas les douze scalaires de \(K\) avec les vingt-quatre coordonnées des fibres \(A_2\), et on ne transforme pas l’action en permutation d’hexades lorsque le relèvement agit également à l’intérieur des fibres.

\section{Portée exacte concernant les histoires}

Soit \(k:X\to H\) le lecteur du registre sur un domaine d’histoires, et soit \(\ell=F\circ k\). Alors

\[
\ell(x)=\ell(y)\quad\Longleftrightarrow\quad k(x)=k(y).
\]

L’écran complet n’ajoute aucune perte par rapport à \(K\), mais il ne change pas non plus les fibres de ce lecteur. Une application \(d:X\to X\) descend en une dynamique sur les valeurs de \(K\) si et seulement si

\[
k(x)=k(y)\Longrightarrow k(d x)=k(d y).
\]

La même condition est nécessaire et suffisante pour descendre sur l’image incidencielle. La preuve consiste à définir \(D(k(x))=k(dx)\) ; la condition équivaut à l’indépendance de la définition à l’égard du représentant. Si ce \(D\) est linéaire et unitaire dans la réalisation déclarée, on applique la section~\ref{x_reciprocidad:nuc05:levantamiento}. S’il est non linéaire, \(F\) continue à le conjuguer comme application, mais ne le transforme pas en unitaire.

Le critère distingue la descente à \(K\) de la dynamique enrichie. La section~\ref{x_reciprocidad:subsec:k-terminal} établit que le lecteur utilise le sous-chemin, l’orientation, la retenue et les incidences ; la section~\ref{x_reciprocidad:exc:limite-inverso} distingue la récupération du mot visible, du support et de la mémoire profonde. Dans \(\ell^2(Z_n)\otimes H\), la section~\ref{x_reciprocidad:nuc05:refinamiento} conserve les étiquettes enrichies et transporte la représentation de la fibre. L’équivalence agit sur le domaine d’histoires préalablement engendré.

Pour une coupe déjà construite, la naturalité du registre dans \cref{x_reciprocidad:prop:k-extension} fixe \(K_N=K_{108}\circ\tau_{N,108}\) : \(F\) conserve exactement cette naturalité. Une suite de nouveaux événements peut mettre à jour \(K\) ; la composition additive ou affine de \cref{x_reciprocidad:sec:k-registro} conserve l’indépendance à l’égard de la partition, et non la constance de la valeur lors de l’ajout d’événements.
